% ==============================================================================
%  Extracting the 0++ glueball mass on the lattice — standalone report
%  (independent of the thesis in tex/; build: pdflatex, run twice for refs)
%
%  Figures are the PNGs produced by the scripts at the repository root;
%  \graphicspath points one level up.
% ==============================================================================
\documentclass[11pt,a4paper]{article}

\usepackage[margin=2.6cm]{geometry}
\usepackage[T1]{fontenc}
\usepackage[utf8]{inputenc}
\usepackage{lmodern}
\usepackage{microtype}
\usepackage{amsmath,amssymb,mathtools}
\usepackage{graphicx}
\usepackage{booktabs}
\usepackage{caption}
\usepackage{xcolor}
\usepackage[colorlinks=true,linkcolor=blue!50!black,citecolor=blue!50!black,urlcolor=blue!50!black]{hyperref}

\graphicspath{{../}}

% ---- shorthand ---------------------------------------------------------------
\newcommand{\Tr}{\operatorname{Tr}}
\newcommand{\dU}{\mathcal{D}U}
\newcommand{\Obar}{\bar{O}}
\newcommand{\meff}{m_{\mathrm{eff}}}
\newcommand{\at}{a_t}
\newcommand{\as}{a_s}
\newcommand{\That}{\hat{T}}

\title{Extracting the $0^{++}$ Glueball Mass on the Lattice:\\
from the Euclidean Path Integral to a Learned Variational Operator}
\author{Francesco Passante}
\date{August 2, 2026}

\begin{document}
\maketitle

\begin{abstract}
Glueballs are bound states of pure gauge fields, and their masses are among the
most natural --- and notoriously noisy --- predictions of lattice gauge theory.
This note develops, from first principles, everything needed to extract the
lightest scalar ($0^{++}$) glueball mass in $SU(2)$ lattice gauge theory:
the Euclidean path integral and the transfer matrix, the spectral decomposition
of connected correlators, the effective mass and its variational bound, the
signal-to-noise catastrophe specific to glueballs, and the three classical
tools that tame it --- link smearing, the variational method (GEVP), and the
anisotropic lattice. We then document an actual extraction on an
anisotropic $SU(2)$ ensemble ($12^3\times24$, $\beta=2.4$, $\xi=3$), where a
multi-level smearing GEVP resolves a clean plateau at $m\,\at\approx0.33$.
Finally, we replace the hand-built operator basis by a gauge-equivariant
neural network (GELT) trained on a Rayleigh-quotient loss whose converged
value \emph{is} the mass. Fed multi-level smeared input channels and
constrained to a single timeslice (so the transfer-matrix bound applies), the
learned operator saturates the variational bound at the classical anchor,
plateaus from $\Delta=1$, and carries significantly less excited-state
contamination than the classical GEVP ($3.9\sigma$ from a correlated-difference
jackknife) while agreeing with it on the mass itself. A from-scratch
replication on an independently sampled ensemble reproduces the result;
the two ensembles combine to a $3.6\sigma$ excess in ground-state overlap.
A closing exploratory section reads the learned operator's gauge-invariant
attention weights out as a measurement, finding its advantage concentrated in
a single local attention head.
\end{abstract}

\tableofcontents

% ==============================================================================
\section{Introduction}
% ==============================================================================

In a pure Yang--Mills theory --- QCD with the quarks removed --- the only
physical states are \emph{glueballs}: color-singlet bound states of gluons,
held together by the same self-interaction that makes the theory confining.
Their spectrum is a parameter-free prediction of the gauge dynamics, and
computing it was one of the first triumphs of lattice gauge theory
\cite{Berg1980,MorningstarPeardon1999}. The lightest state carries the quantum
numbers of the vacuum, $J^{PC}=0^{++}$, and it is the target of everything
that follows.

Extracting a glueball mass is conceptually simple --- measure a Euclidean
correlator, read off its exponential decay --- and practically hard: the
signal decays exponentially while the noise does not, so the mass must be read
at very small time separations, which is only possible if the measuring
operator overlaps strongly onto the ground state. The entire craft of glueball
spectroscopy is the construction of such operators.

This note has two halves. Sections~\ref{sec:pathintegral}--\ref{sec:classical}
build the standard theory and practice from the ground up, assuming only
basic familiarity with lattice gauge theory, and document a classical
measurement of the $SU(2)$ $0^{++}$ mass on an anisotropic lattice.
Sections~\ref{sec:learned}--\ref{sec:runs} then reformulate the operator
construction as a machine-learning problem: since the optimal operator is the
maximizer of a Rayleigh quotient of the transfer matrix, a gauge-equivariant
network trained to minimize $-C(1)/C(0)$ \emph{is} a variational glueball
operator, and its converged loss \emph{is} (minus the exponential of) the
mass. The final result is a learned operator that variationally subsumes the
hand-built smearing basis it was validated against.

% ==============================================================================
\section{From the path integral to the lattice}
\label{sec:pathintegral}
% ==============================================================================

\subsection{The Euclidean path integral}

Continuum expectation values of a quantum field theory can be written, after
Wick rotation $t\to -i\tau$, as statistical averages with a real, positive
weight,
\begin{equation}
  \langle O \rangle
  \;=\;
  \frac{1}{Z}\int \mathcal{D}A\; O[A]\, e^{-S_E[A]},
  \qquad
  Z=\int \mathcal{D}A\; e^{-S_E[A]},
  \label{eq:pathintegral}
\end{equation}
where $S_E$ is the Euclidean action. Two things make this formulation the
foundation of lattice gauge theory. First, $e^{-S_E}$ is a probability weight
(up to normalization), so \eqref{eq:pathintegral} can be estimated by Monte
Carlo importance sampling. Second --- and this is the part that matters for
spectroscopy --- Euclidean time evolution is generated by $e^{-\tau\hat H}$
rather than $e^{-it\hat H}$: excited states are \emph{exponentially
suppressed} rather than oscillating, so the large-$\tau$ behaviour of
correlation functions isolates the lowest-energy state in a given channel.
Masses are read off exponential decays. Section~\ref{sec:correlators} makes
this precise via the transfer matrix.

\subsection{Gauge fields on a lattice: links, plaquettes, Wilson action}

Wilson's discretization \cite{Wilson1974} replaces spacetime by a hypercubic
lattice $\Lambda$ with spacing $a$, and the gauge potential $A_\mu(x)$ by
group-valued \emph{link variables}
\begin{equation}
  U_\mu(x) \;=\; \mathcal{P}\exp\!\Big( i g \!\int_x^{x+a\hat\mu}\! A\cdot dx \Big)
  \;\in\; SU(N),
\end{equation}
the parallel transporters along the elementary edges of the lattice. A gauge
transformation is a group element $\Omega(x)$ at every site, acting as
\begin{equation}
  U_\mu(x) \;\longrightarrow\; \Omega(x)\, U_\mu(x)\, \Omega^\dagger(x+a\hat\mu).
  \label{eq:gaugetrafo}
\end{equation}
Any product of links along a \emph{closed} loop therefore transforms by
conjugation at its base point, and its trace is exactly gauge invariant. The
smallest such loop is the \emph{plaquette},
\begin{equation}
  P_{\mu\nu}(x)
  \;=\;
  U_\mu(x)\, U_\nu(x+a\hat\mu)\, U_\mu^\dagger(x+a\hat\nu)\, U_\nu^\dagger(x),
\end{equation}
and the Wilson action is the sum of plaquettes,
\begin{equation}
  S[U] \;=\; \beta \sum_{x,\,\mu<\nu}
  \Big( 1 - \tfrac{1}{N}\,\mathrm{Re}\,\Tr P_{\mu\nu}(x) \Big),
  \qquad
  \beta = \frac{2N}{g^2},
  \label{eq:wilson}
\end{equation}
which reproduces $\tfrac14 F_{\mu\nu}^aF_{\mu\nu}^a$ in the classical
continuum limit $a\to0$. The bare coupling $\beta$ is the only parameter;
the lattice spacing $a(\beta)$ is not an input but an output, set by the
dynamics (asymptotic freedom drives $a\to0$ as $\beta\to\infty$). All results
below are quoted in lattice units, i.e.\ as dimensionless products such as
$m\,a$.

\subsection{Expectation values by Monte Carlo}

The lattice path integral is an ordinary (huge-dimensional) integral over the
compact group manifold, with Haar measure on every link:
\begin{equation}
  \langle O \rangle
  \;=\;
  \frac{1}{Z}\int \prod_{x,\mu} dU_\mu(x)\; O[U]\, e^{-S[U]} .
\end{equation}
A Markov-chain Monte Carlo algorithm generates an ensemble of configurations
$\{U^{(1)},\dots,U^{(N)}\}$ distributed as $e^{-S}$, and
$\langle O\rangle \approx \frac1N \sum_i O[U^{(i)}]$ with an error
$\propto 1/\sqrt{N_{\rm eff}}$, where $N_{\rm eff}$ accounts for
autocorrelation along the chain (Section~\ref{sec:ensemble}). Everything in
this note runs in the gauge group $SU(2)$, which is cheap, confining, and
admits exact single-link updates.

% ==============================================================================
\section{Masses from Euclidean correlators}
\label{sec:correlators}
% ==============================================================================

\subsection{The transfer matrix}

The Euclidean lattice theory on a lattice with $N_t$ timeslices and periodic
temporal boundary conditions is equivalent to a quantum system with
Hamiltonian $\hat H$ acting on a Hilbert space of spatial-field states: the
partition function factorizes as
\begin{equation}
  Z \;=\; \Tr\, \That^{\,N_t},
  \qquad
  \That \;=\; e^{-\at \hat H},
\end{equation}
where $\That$, the \emph{transfer matrix}, propagates the system by one
timeslice ($\at$ is the temporal lattice spacing). For the Wilson action
$\That$ is self-adjoint and positive (\emph{reflection positivity}
\cite{OsterwalderSeiler1978}), which is what makes everything below a theorem
rather than a heuristic: $\That$ has a complete set of eigenstates $|n\rangle$
with real eigenvalues $e^{-\at E_n}$, $E_0 \le E_1 \le \dots$, and we measure
energies relative to the vacuum, $E_0 = 0$.

An operator $O(\vec x, t)$ built from the \emph{spatial} links of timeslice
$t$ only corresponds to a Schr\"odinger-picture operator $\hat O$ on this
Hilbert space. This locality-in-time condition looks like a technicality; it
will return twice with teeth --- once for smearing
(Section~\ref{sec:smearing}) and once, decisively, for the neural network
(Section~\ref{sec:constraint}).

\subsection{Spectral decomposition of the correlator}

Insert complete sets of transfer-matrix eigenstates into a two-point function
at temporal separation $\Delta$ (in units of $\at$):
\begin{equation}
  \big\langle O(t{+}\Delta)\, O(t) \big\rangle
  \;=\;
  \frac{1}{Z}\Tr\!\big[ \That^{\,N_t-\Delta}\, \hat O\, \That^{\,\Delta}\, \hat O \big]
  \;\xrightarrow[N_t\gg\Delta]{}\;
  \sum_{n} \big|\langle 0|\hat O|n\rangle\big|^2\, e^{-E_n \Delta\,\at}.
  \label{eq:spectral}
\end{equation}
Every term is positive, and the decay rates are the energies of the states
the operator couples to. At large $\Delta$ the lightest such state dominates:
its energy is read off the exponential tail. (On a finite periodic lattice
the exponential is replaced by
$e^{-E_n\Delta \at}+e^{-E_n(N_t-\Delta)\at}\propto\cosh$, symmetric about
$\Delta=N_t/2$; this matters only near the midpoint.)

\subsection{Zero momentum, vacuum subtraction, effective mass}

Three standard refinements turn \eqref{eq:spectral} into a mass measurement.

\paragraph{Zero-momentum projection.} Summing the operator over a spatial
slice,
\begin{equation}
  \Obar(t) \;=\; \sum_{\vec x} O(\vec x, t),
\end{equation}
projects onto momentum $\vec p=0$, so the energies in the decomposition are
\emph{masses}, $E_n = m_n$, with no kinetic contamination. It also averages
over $L^3$ sites, improving statistics.

\paragraph{Vacuum subtraction.} The $0^{++}$ channel shares the quantum
numbers of the vacuum, so $\langle\Obar\rangle\neq0$ and the $n=0$ term of
\eqref{eq:spectral} is a $\Delta$-independent constant
$|\langle0|\hat O|0\rangle|^2$ that would bury the signal. One therefore
measures the \emph{connected} correlator,
\begin{equation}
  C(\Delta)
  \;=\;
  \big\langle \Obar(t_0{+}\Delta)\,\Obar(t_0) \big\rangle
  - \big\langle \Obar \big\rangle^2
  \;=\;
  \sum_{n>0} \big|\langle 0|\hat{\bar O}|n\rangle\big|^2\, e^{-m_n \Delta\,\at},
  \label{eq:connected}
\end{equation}
averaged over all time origins $t_0$ (time-translation invariance buys another
factor $N_t$ of statistics). A structural gift of the $0^{++}$ channel: since
it \emph{is} the vacuum channel, after subtraction the lightest surviving
state is automatically the scalar glueball --- no quantum-number projection
(spin, parity, charge conjugation on the cubic group) is needed for the
ground state.

\paragraph{Effective mass.} Define
\begin{equation}
  \meff(\Delta) \;=\; \log\frac{C(\Delta)}{C(\Delta+1)} .
  \label{eq:meff}
\end{equation}
Because every term in \eqref{eq:connected} is positive,
$\meff(\Delta)\ge m_G$ for all $\Delta$ and $\meff(\Delta)\to m_G\,\at$
monotonically from above as $\Delta\to\infty$: the correlator is a
\emph{variational upper bound} machine. In practice one looks for the
\emph{plateau} --- the range of $\Delta$ where $\meff$ has stopped descending
(excited states have decayed) but the error bars have not yet exploded. The
whole game is to make that window exist.

% ==============================================================================
\section{Why the glueball is hard: signal-to-noise}
\label{sec:snr}
% ==============================================================================

For most hadrons the plateau window is wide. For glueballs it is a crack in
the wall, for a reason first made sharp by Lepage \cite{Lepage1989}. The Monte
Carlo variance of the correlator estimate is itself a correlator,
\begin{equation}
  \mathrm{Var}\big[ C(\Delta) \big]
  \;\sim\;
  \big\langle \Obar^2(t_0{+}\Delta)\, \Obar^2(t_0) \big\rangle - C(\Delta)^2 ,
\end{equation}
and the operator $\Obar^2$ couples to the \emph{vacuum}: the variance
approaches a $\Delta$-independent constant. The signal decays as
$e^{-m_G\Delta \at}$ while the noise does not, so
\begin{equation}
  \frac{\text{signal}}{\text{noise}}
  \;\sim\;
  \sqrt{N}\; e^{-m_G \Delta\, \at} .
  \label{eq:lepage}
\end{equation}
The scalar glueball is \emph{heavy} --- $m_G\, a$ is order one on typical
lattices --- so the signal is gone within two or three timeslices, and by
\eqref{eq:lepage} doubling the statistics buys only $\log\sqrt2/(m_G \at)$
extra slices. More configurations are nearly useless; the plateau must be
dragged to \emph{small} $\Delta$, where the noise is still small. Since
$\meff(\Delta)$ starts high and descends as excited states die off, dragging
the plateau to small $\Delta$ means building an operator with large
\emph{ground-state overlap}: $|\langle 0|\hat{\bar O}|1\rangle|^2$ dominant in
\eqref{eq:connected} already at $\Delta=1$. Everything in the next section is
a technique for manufacturing overlap.

% ==============================================================================
\section{The classical toolkit}
\label{sec:toolkit}
% ==============================================================================

\subsection{Smeared operators}
\label{sec:smearing}

The simplest scalar glueball operator is the spatial plaquette itself, summed
over orientations to make a rotational scalar. It is a terrible interpolating
operator: a single \emph{thin-link} loop of extent $a$ probes UV fluctuations,
not the physical glueball, whose wavefunction has a size of order the
confinement scale. Its $\meff$ starts far above $m_G$ and descends too slowly
to plateau before the noise wall of \eqref{eq:lepage}.

The fix is \emph{link smearing}: replace each spatial link by a
gauge-covariant local average of itself and its neighbouring
\emph{staples} (the three-link detours connecting the same endpoints). One
APE step \cite{APE1987} is
\begin{equation}
  U_j(x)\;\longrightarrow\;
  \mathrm{Proj}_{SU(N)}\Big[
    (1-\alpha)\, U_j(x)
    + \frac{\alpha}{n_{\rm st}} \sum_{\text{spatial staples}} S_j(x)
  \Big],
  \label{eq:ape}
\end{equation}
iterated $n$ times; operators built from smeared links are spatially extended
over a radius that grows with $n$, matching the physical size of the state and
suppressing the UV contamination. Two rules are absolute:
\begin{itemize}
  \item \textbf{Smear spatial links only, never in time.} Temporal smearing
  makes the operator depend on several timeslices at once, destroying the
  transfer-matrix interpretation of \eqref{eq:spectral} --- the ``operator on
  timeslice $t$'' premise --- and with it the variational bound.
  \item \textbf{Smearing is not cooling.} Cooling/gradient flow drives the
  \emph{configuration} toward classical solutions and changes the measured
  physics; operator smearing changes only the \emph{operator}, leaving the
  ensemble untouched. Both are built from staples; they answer different
  questions.
\end{itemize}
There is no single best smearing level --- too little leaves UV
contamination, too much washes the operator out --- which motivates using
\emph{several} levels at once:

\subsection{The variational method (GEVP)}
\label{sec:gevp}

Given a basis of $n_{\rm op}$ operators $\Obar_i$ (here: the scalar operator
at APE levels $0,2,4,6$), measure the full connected correlator
\emph{matrix}
\begin{equation}
  C_{ij}(\Delta) \;=\;
  \big\langle \Obar_i(t_0{+}\Delta)\, \Obar_j(t_0) \big\rangle
  - \langle\Obar_i\rangle \langle\Obar_j\rangle ,
\end{equation}
and solve the \emph{generalized eigenvalue problem} (GEVP)
\cite{MichaelTeasdale1983,LuscherWolff1990,Blossier2009}
\begin{equation}
  C(\Delta)\, v_n \;=\; \lambda_n(\Delta, t_0)\; C(t_0)\, v_n ,
  \label{eq:gevp}
\end{equation}
at a fixed reference time $t_0$ (here $t_0=1$). The eigenvector $v_n$ defines
the linear combination $\sum_i (v_n)_i \Obar_i$ that optimally isolates state
$n$ within the span of the basis, and the eigenvalues decay with the
\emph{individual} state energies,
\begin{equation}
  \lambda_n(\Delta,t_0) \;=\; e^{-m_n (\Delta-t_0)\,\at}
  \big( 1 + \mathcal{O}(e^{-\delta m\, \Delta \,\at}) \big),
\end{equation}
with corrections controlled by the gap $\delta m$ to the first state
\emph{outside} the basis span --- parametrically smaller than the
single-operator contamination. Effective masses are then read from the
eigenvalues, $\meff^{(n)}(\Delta)=\log[\lambda_n(\Delta)/\lambda_n(\Delta+1)]$,
for $\Delta\ge t_0$. This ``multi-level smearing basis + GEVP'' construction
is the Morningstar--Peardon strategy \cite{MorningstarPeardon1999} that
underlies every modern glueball spectrum.

One numerical remark from practice: at low statistics the estimated $C(t_0)$
can be indefinite, so a naive Cholesky-based solver fails. The implementation
used here whitens with the eigendecomposition of $C(t_0)$ and floors
near-zero eigenvalues --- adequate for four well-separated smearing levels,
though for larger bases the flooring should be replaced by truncation
(projecting the noisy directions out rather than blowing them up).

\subsection{Anisotropic lattices}
\label{sec:aniso}

Smearing and the GEVP fix the \emph{overlap}; they cannot manufacture
timeslices. If $m_G\, \at \sim 1$, the correlator falls by $e^{-1}$ per slice
and even a perfect operator has only $2$--$3$ usable points --- too few to
demonstrate a plateau. The field's standard cure
\cite{MorningstarPeardon1999} is the \emph{anisotropic lattice}: a temporal
spacing finer than the spatial one, $\at = \as/\xi$ with anisotropy
$\xi>1$, so the same physical decay is sampled by $\xi$ times more slices.
At tree level this is a reweighting of the Wilson action by plaquette
orientation,
\begin{equation}
  S \;=\; \beta_s \!\!\sum_{\text{spatial } P}\!\!
        \Big(1-\tfrac1N \mathrm{Re}\Tr P\Big)
      + \beta_t \!\!\sum_{\text{temporal } P}\!\!
        \Big(1-\tfrac1N \mathrm{Re}\Tr P\Big),
  \qquad
  \beta_s = \frac{\beta}{\xi},\quad \beta_t = \beta\,\xi .
  \label{eq:aniso}
\end{equation}
The masses extracted from temporal decays come out in units of $\at$; the
conversion $m\,\as = \xi\, m\,\at$ uses the \emph{renormalized} anisotropy
$\xi_R$, which quantum corrections shift away from the bare $\xi$ (we measure
the shift but do not tune it away; see Section~\ref{sec:caveat}).

% ==============================================================================
\section{Ensemble generation and error analysis}
\label{sec:ensemble}
% ==============================================================================

\subsection{Heat-bath and overrelaxation for $SU(2)$}

Spectroscopy needs many well-decorrelated configurations, and the workhorse
Metropolis algorithm decorrelates slowly (each accepted move is a small random
step). $SU(2)$ admits two exact, tuning-free local updates that together beat
this critical slowing down:

\paragraph{Heat-bath \cite{Creutz1980,KennedyPendleton1985}.} The action is
linear in any single link: $S \supset -\tfrac{\beta}{N}\mathrm{Re}\Tr[U_\mu(x) A^\dagger]$,
where $A$ is the sum of the staples around that link (with the $\xi^{\pm1}$
weights of \eqref{eq:aniso} folded in). For $SU(2)$ a sum of group elements is
proportional to a group element, $A = k\,V$ with $V\in SU(2)$, so the
single-link conditional distribution is known in closed form and can be
sampled \emph{exactly} --- the new link is completely independent of the old
one.

\paragraph{Overrelaxation \cite{Creutz1987}.} The reflection
$U \to V^\dagger U^\dagger V^\dagger$ preserves the action exactly
(microcanonical) while moving as far as possible across the constant-action
surface --- pure decorrelation at negligible cost. The production sweep used
below is one heat-bath sweep followed by $n_{\rm OR}=4$ overrelaxation sweeps,
in a checkerboard (even/odd) pattern so that whole sublattices update in
parallel.

\subsection{Autocorrelation}

Successive configurations of the chain are correlated; the effective sample
size is $N_{\rm eff} = N / (2\tau_{\rm int})$, where the integrated
autocorrelation time $\tau_{\rm int}$ of the \emph{measured observable}
(estimated with Madras--Sokal automatic windowing \cite{MadrasSokal1988}) sets
the number of sweeps $n_{\rm skip}\gtrsim 2\tau_{\rm int}$ to discard between
saved configurations. Crucially, $\tau_{\rm int}$ must be measured on the
smeared glueball operator, not the plaquette --- extended observables
decorrelate more slowly.

\subsection{Jackknife errors}

All quoted errors are leave-one-out jackknife errors over configurations:
recompute the full analysis (correlator, GEVP, $\meff$) on $N$ subsamples each
omitting one configuration; the spread of the results estimates the error of
the full-sample result, propagating correctly through every nonlinearity
(ratios, logs, eigenvalues). When residual autocorrelation is suspected, the
omitted unit is a contiguous \emph{block} of configurations rather than one
(``blocked jackknife''), which keeps the independence assumption honest. Two
further uses appear below: neighbouring-$\Delta$ points of a jackknifed
$\meff$ curve are strongly correlated, so ``the error bars overlap'' can
overstate consistency; and when comparing two estimates measured on the
\emph{same} configurations, the jackknife of their \emph{difference} cancels
the correlated fluctuations and is the honest significance test
(Section~\ref{sec:run5}).

% ==============================================================================
\section{The classical measurement, run by run}
\label{sec:classical}
% ==============================================================================

All ensembles are $SU(2)$, sampled with heat-bath + overrelaxation. The
operator is the orientation-summed spatial $1\times1$ Wilson loop
(the spatial plaquette scalar), at APE levels $0,2,4,6$; the analysis is the
connected correlator \eqref{eq:connected}, $\meff$ \eqref{eq:meff}, and the
GEVP \eqref{eq:gevp} with $t_0=1$, all errors jackknifed.

\subsection{Runs 0--1: the isotropic lattice fails}

On an isotropic $12^4$ lattice at $\beta=2.4$ with $N=2000$ configurations,
the single smeared operator shows only a weak $m\,a\approx0.8$ at
$\Delta=1$--$2$ and slides into noise (negative $\meff$, huge bars) by
$\Delta\approx3$ --- the textbook heavy-$0^{++}$ failure of
Section~\ref{sec:snr}, verbatim. Adding the full multi-level GEVP (Run~1)
produced a flatter, lower ground state but still \emph{no plateau}: with the
signal dead by $\Delta\approx3$ there is nothing late-$\Delta$ for any basis
to exploit. The diagnosis, important because it redirects the effort: the
\emph{operator basis was not the bottleneck; the lattice was.} No amount of
overlap engineering survives $m_G\,a\sim1$ on an isotropic lattice.

\subsection{Run 2: validating the anisotropic implementation}

Before trusting physics from the anisotropic action \eqref{eq:aniso}, three
checks (Figure~\ref{fig:aniso}):
\begin{itemize}
  \item \textbf{$\xi=1$ regression:} the anisotropic code path at $\xi=1$
  reproduces the exactly known 2D $SU(2)$ mean plaquette
  $I_2(\beta)/I_1(\beta)$: measured $0.4329\pm0.0017$ vs.\ exact $0.4331$
  ($0.1\sigma$) --- the refactor changed nothing where it must not.
  \item \textbf{Plaquette splitting:} in 4D, $\langle P_{st}\rangle$ rises and
  $\langle P_{ss}\rangle$ falls with $\xi$ (since $\beta_t>\beta_s$),
  coinciding at $\xi=1$ --- the anisotropy acts in the intended direction.
  \item \textbf{Renormalized anisotropy:} a Creutz-ratio-based estimate gives
  $\xi_R \approx 3.32$ at bare $\xi=3$ --- the expected tree-level vs.\
  renormalized mismatch, made visible rather than tuned away.
\end{itemize}

\begin{figure}[tbp]
  \centering
  \includegraphics[width=\textwidth]{anisotropy_validation.png}
  \caption{Validation of the anisotropic $SU(2)$ action. Left: the $\xi=1$
  code path reproduces the exact 2D mean plaquette $I_2(\beta)/I_1(\beta)$ to
  $0.1\sigma$. Right: at $\beta=2.0$ in 4D the plaquette splits with the bare
  anisotropy --- temporal plaquettes grow, spatial ones shrink --- and a
  Creutz-ratio estimate gives a renormalized $\xi_R\approx3.32$ at bare
  $\xi=3$.}
  \label{fig:aniso}
\end{figure}

\subsection{Run 3: the anchor, $m\,\at \approx 0.33$}

The production ensemble is $12^3\times24$ at $\beta=2.4$, bare $\xi=3.0$,
$N=2000$ configurations (heat-bath + overrelaxation, acceptance~1.00). The
anisotropy transforms the measurement (Figure~\ref{fig:classical}):
$C(\Delta)$ now decays cleanly over $\sim$10--12 timeslices before the
periodic turnaround at $\Delta=N_t/2=12$, instead of dying at $\Delta\approx3$.
The GEVP ground state (levels $[0,2,4,6]$, $t_0=1$) plateaus with tight
errors:
\begin{equation}
  \meff(1) = 0.365\pm0.008, \quad
  \meff(2) = 0.333\pm0.011
  \qquad (\text{units of } \at),
\end{equation}
a small descent then flattening --- the expected approach from above of the
variational bound. Because two points that differ by $\sim2\sigma$ are a
descent, not yet a plateau, the anchor was later strengthened with
$\Delta=3$--$4$ points on the held-out test split of the same ensemble
(Section~\ref{sec:run4}): $\meff(3)=0.333\pm0.036$ confirms
\begin{equation}
  \boxed{\; m\,\at \approx 0.33 \;}
\end{equation}
as the classical ground-truth anchor for everything that follows.

\begin{figure}[tbp]
  \centering
  \includegraphics[width=\textwidth]{glueball_validation.png}
  \caption{The classical baseline on the anisotropic ensemble
  ($12^3\times24$, $\beta=2.4$, $\xi=3$, $N=2000$). Top left: the analysis
  code recovers a known mass from a synthetic single-exponential correlator.
  Top right: the smeared operator's expectation value grows monotonically
  with APE steps (smearing sanity check). Bottom left: the connected
  correlator $C(\Delta)$ for thin and $6\times$-smeared operators decays over
  $\sim$10--12 slices, symmetric about $\Delta=N_t/2$ (periodicity). Bottom
  right: effective masses --- the multi-level GEVP ground state (green)
  plateaus where the single operators do not.}
  \label{fig:classical}
\end{figure}

\subsection{What the number does and does not mean}
\label{sec:caveat}

$m\,\at\approx0.33$ is the trustworthy, method-validating result. The
conversion $m\,\as = \xi\, m\,\at \approx 1.0$ is \emph{not} continuum
physics: at $\beta_s=\beta/\xi=0.8$ the spatial lattice is deep in strong
coupling (coarse $\as$), so $m\,\as$ carries large discretization artifacts,
and the conversion moreover uses the bare $\xi=3$ where Run~2 measured
$\xi_R\approx3.32$. A publishable $m_G$ would require tuned anisotropy
(raising $\beta$ with $\xi$ to keep $\beta_s$ in the scaling window) and a
continuum extrapolation --- deliberately left as future work. The claim
established here is the one the rest of this note needs: \emph{the pipeline
resolves a plateau}, i.e.\ a classical ground truth exists on this ensemble.

% ==============================================================================
\section{A neural network as a variational operator}
\label{sec:learned}
% ==============================================================================

Everything so far is a search for the operator with maximal ground-state
overlap, conducted by hand: choose loop shapes, choose smearing levels, let a
GEVP take one global linear combination. The observation driving the second
half of this note is that this search is an \emph{optimization problem with a
differentiable objective}, and the objective is exactly a physics quantity.

\subsection{The Rayleigh-quotient loss}
\label{sec:rayleigh}

Let $O_\theta(\vec x, t)$ be \emph{any} gauge-invariant scalar field built
from the spatial links of timeslice $t$ (with parameters $\theta$), and
$\Obar_\theta(t)$ its zero-momentum projection. Write
$|\psi\rangle = (\hat{\bar O}_\theta - \langle\Obar_\theta\rangle)\,|0\rangle$
for the vacuum-subtracted state it creates. Then by construction
$C(\Delta) = \langle\psi|\That^\Delta|\psi\rangle$ and $\langle\psi|0\rangle=0$,
so the correlator ratio
\begin{equation}
  R(\theta) \;=\; \frac{C(1)}{C(0)}
  \;=\; \frac{\langle\psi|\That|\psi\rangle}{\langle\psi|\psi\rangle}
  \label{eq:rayleighquotient}
\end{equation}
is a \emph{Rayleigh quotient of the transfer matrix} restricted to the
orthogonal complement of the vacuum. Since $\That$ is positive and
self-adjoint (reflection positivity again), the supremum of
\eqref{eq:rayleighquotient} over all operators is the largest non-vacuum
eigenvalue,
\begin{equation}
  R(\theta) \;\le\; e^{-m_G\,\at},
\end{equation}
with equality exactly when $|\psi\rangle$ is the glueball ground state.
Therefore:
\begin{quote}
\textbf{Train a network on the loss $\mathcal{L}(\theta) = -\,C(1)/C(0)$.}
The loss is a rigorous variational bound; at convergence
$m_G\,\at = -\log R$; the converged loss value \emph{is} the glueball mass,
and the trained network \emph{is} the optimal interpolating operator.
\end{quote}
No labels, no supervision --- the training signal is the same Monte Carlo
correlator the classical analysis measures. Two structural gifts: the
quotient is invariant under $O\to\lambda O$, so the output needs no
normalization (this gift has a sting; Section~\ref{sec:pitfalls}); and, as in
Section~\ref{sec:correlators}, the $0^{++}$ channel needs no quantum-number
projection for the ground state. In practice the loss used below averages two
ratios, $\mathcal{L} = -\tfrac12\,[\,C(1)/C(0) + C(2)/C(0)\,]$, whose floor at
the anchor mass is $-\tfrac12(e^{-m\at}+e^{-2m\at}) \approx -0.62$ at
$m\,\at=0.33$ --- a number to remember.

\subsection{The GELT architecture in brief}

The operator family $O_\theta$ is GELT, a gauge-equivariant lattice
transformer in the spirit of the L-CNN framework \cite{Favoni2020}. Only its
contract matters here:
\begin{itemize}
  \item \textbf{Input:} untraced plaquettes $W(x)$, which transform
  \emph{adjointly}, $W(x)\to\Omega(x)W(x)\Omega^\dagger(x)$, plus parallel
  transporters $T_{\Delta x}(x)$ to every site in an $L^1$-ball of radius $R$,
  averaged over all shortest lattice paths.
  \item \textbf{Attention:} scores are gauge \emph{invariant}
  ($\mathrm{Re}\Tr[Q^\dagger(x)\,\tilde K(x{+}\Delta x)]$, with $K$
  transported to $x$), so the attention map itself is a gauge-invariant,
  physically interpretable object.
  \item \textbf{Values:} the value path is matrix-\emph{bilinear},
  $\sum \alpha\, Q_v^\dagger \tilde V$ --- each layer multiplies transported
  loop matrices, so depth composes ever-larger Wilson loops, the same
  universality mechanism as L-CNN's bilinear layers.
  \item \textbf{Output:} a trace and per-site MLP produce a gauge-invariant
  scalar field $O_\theta(x)$ per site --- precisely an operator density.
\end{itemize}
Equivariance throughout guarantees the output is exactly gauge invariant, so
$O_\theta$ is a legitimate physical operator for any $\theta$: the
variational bound of Section~\ref{sec:rayleigh} holds at every step of
training, not just at convergence. Conceptually, GELT is a \emph{learned,
content-dependent smearing}: transport-averaging plus stacked attention is a
gauge-covariant multi-scale aggregation, and the Rayleigh loss tunes it for
ground-state overlap directly --- the thing hand-tuned APE schedules
approximate.

\subsection{The single-timeslice constraint (where the theory bites)}
\label{sec:constraint}

The derivation of \eqref{eq:rayleighquotient} required $\Obar(t)$ to be a
functional of the fields on \emph{timeslice $t$ only} --- the same rule that
forbids temporal smearing (Section~\ref{sec:smearing}), now applied to the
network itself. A 4D network violates it doubly: temporal plaquettes among
its input channels, and a receptive field (transport radius $\times$ depth)
that leaks across timeslices. The failure is not a small correction; it makes
the loss \emph{gameable}. If $O(t)$ may depend on neighbouring slices, the
optimizer's cheapest route to $C(1)/C(0)\to1$ is an output that varies slowly
in $t$ --- in the limit, a per-configuration constant gives
$C(\Delta)=C(0)$ for all $\Delta$ and a fake mass of zero. The network would
``beat'' every classical operator for a purely unphysical reason.

The fix defines the setup: \textbf{run GELT as a three-dimensional network on
the spatial links of one timeslice}, batched over
(configuration $\times$ timeslice), with per-timeslice outputs reshaped to
$\Obar(t)$ for the correlator. Within its slice the operator keeps its full
spatial loop-building power --- exactly the domain the classical operator
uses. The absence of temporal loop content is not a limitation; it is the
definition of a valid interpolating operator. (The 3D restriction also
collapses the transport memory cost by $\sim70\times$, making on-the-fly
transport per batch cheap --- the constraint that protects the physics also
pays the compute bill.)

\subsection{Making the loss trainable}
\label{sec:pitfalls}

Four practical pathologies, each discovered the hard way:

\paragraph{Zero gradient at initialization.} The network's readout was
zero-initialized (a common stabilization). But $C(0)$ and $C(1)$ are both
quadratic in the output, so at $O\equiv0$ the gradient of $-C(1)/(C(0)+\epsilon)$
is \emph{exactly} zero --- training never leaves the saddle. The readout must
be initialized nonzero.

\paragraph{Denominators.} An early loss chained ratios
$C(\Delta)/C(\Delta{-}1)$, putting the small, signed, noisy $C(1)$ in a
denominator; clamping kept the values finite but the \emph{gradients} were
garbage in direction, and the operator scale was pumped to $C(0)\sim5\times10^8$.
Every ratio now has the large, positive $C(0)$ as its denominator, and the
correlator arithmetic runs in float64.

\paragraph{The scale-flat direction.} Scale invariance of the quotient means
$\theta$-space has an exactly flat direction $O\to\lambda O$. With smooth,
site-coherent inputs (Section~\ref{sec:run5}) the bilinear value path squares
its gain per layer, and the optimizer drifted along the flat direction to
$C(0)\sim10^{73}$ --- float32 overflow, \texttt{NaN} validation loss, and
(insidiously) no checkpoint ever saved, because \texttt{NaN} never improves
the best-so-far. The cure is a variationally neutral \emph{scale pin}: adding
$(\log C(0))^2$ to the loss penalizes only the flat direction (it is minimized
at $C(0)=1$ and does not touch the ratio), removing the runaway without
biasing the physics.

\paragraph{Estimator bias and overfitting.} The loss is a ratio of
\emph{batch-estimated} means (including the vacuum subtraction, itself a
noisy batch estimate), hence biased on small batches --- so batches are made
as large as memory allows and the final numbers never come from the training
estimator. More fundamentally, the variational bound holds in
\emph{expectation}, not on a finite sample: a network can overfit the noise
of its training correlator and spuriously ``violate'' the bound there. The
protocol is therefore a hard split (70\% train / 10\% validation / 20\% test,
contiguous in Monte Carlo time), checkpoint selection on validation, and every
reported mass from a blocked jackknife on the untouched test configurations
only.

% ==============================================================================
\section{Training runs: from a failed operator to one that wins}
\label{sec:runs}
% ==============================================================================

All runs train on the cached Run-3 ensemble ($12^3\times24$, $\beta=2.4$,
$\xi=3$, $N=2000$; 1400 train / 200 validation / 400 test configurations),
with a small per-timeslice 3D GELT ($R=2$, 4 layers, 2 heads,
$d_{\rm qkv}=6$, $d_{\rm model}=16$; $\sim$5k parameters), AdamW with cosine
schedule, batches of 6 configurations ($=144$ timeslices), and gradient
checkpointing on a V100. The classical GEVP of Section~\ref{sec:classical},
re-evaluated on the same test split with the same blocked jackknife, is the
baseline throughout.

\subsection{Run 4 --- thin input: a genuine operator that loses on overlap}
\label{sec:run4}

Fed thin spatial plaquettes, GELT trains stably to a best validation loss of
$-0.433$ --- well short of the $\approx-0.62$ variational floor --- with a
clear overfitting signature (validation bottoms out near epoch 12 while the
training loss keeps falling; Figure~\ref{fig:run4}, left). On test:

\begin{table}[htbp]
  \centering
  \begin{tabular}{lcccc}
    \toprule
    operator & $\meff(1)$ & $\meff(2)$ & $\meff(3)$ & $\meff(4)$ \\
    \midrule
    GELT (thin input)   & $0.527\pm0.028$ & $0.441\pm0.042$ & $0.310\pm0.045$ & $\sim0.5$, large bars \\
    classical GEVP      & $0.398\pm0.016$ & $0.357\pm0.025$ & $0.333\pm0.036$ & $0.284\pm0.050$ \\
    \bottomrule
  \end{tabular}
  \caption{Run 4 (thin-input GELT) vs.\ the classical GEVP, blocked jackknife
  on the 400 test configurations. Units of $\at$.}
  \label{tab:run4}
\end{table}

The learned operator is \emph{genuine} --- its $\meff$ descends
$0.705\to0.527\to0.441\to0.310$, the $\Delta=3$ point consistent with the
anchor, and it clearly beats the thin classical operator
($\approx0.8$ at $\Delta=1$) with no hand-built smearing. But it carries
excited-state contamination the smeared GEVP basis has projected out, and no
convincing plateau of its own.

The diagnosis is a \emph{depth budget}. Four bilinear layers compose loops of
degree $\le16$ in the input plaquettes; six iterated APE steps contain
branched staple content of degree $\sim3^6$ in the links within the same
spatial radius. The network cannot re-derive an iterated smearing schedule at
affordable depth. This diagnosis was made quantitative by an eval-only
diagnostic (Run 4b): the correlation of the learned $\Obar$ with the APE
ladder on test fluctuations peaks at APE$\times2$
($|r| = 0.68,\ \mathbf{0.87},\ 0.71,\ 0.59$ for levels $0,2,4,6$) --- the
depth-4 network learned about two APE steps' worth of staple content.
Run 4b also appended the learned operator to the classical basis as a fifth
GEVP operator: the enlarged basis is \emph{identical} to the classical one
within small fractions of the error bars (Figure~\ref{fig:run4}, the
overlapping curves). The $\sim25\%$ of the learned operator's variance that
is linearly independent of the ladder carries no additional ground-state
overlap --- the thin-input operator is variationally redundant. This null is
the clean baseline for what follows: any future separation of the enlarged
GEVP from the classical curve must come from what the network computes
\emph{on top of} smearing, not from smearing itself.

\begin{figure}[tbp]
  \centering
  \includegraphics[width=\textwidth]{glueball_gelt_from_raw_plaquettes.png}
  \caption{Run 4: GELT trained from thin spatial plaquettes. Left: Rayleigh
  training curves; validation (orange) bottoms out near epoch 12 and turns up
  --- overfitting of the empirical correlator --- while the best value
  $-0.433$ stays far from the $\approx-0.62$ variational floor. Right: test
  effective masses. The learned operator (green) is genuine (descending
  toward the anchor, far better than the thin operator) but sits above the
  classical GEVP (red) at small $\Delta$: an overlap gap.}
  \label{fig:run4}
\end{figure}

\subsection{Run 5 --- smeared input channels: the physics fix}
\label{sec:run5}

If depth cannot rebuild iterated smearing, give the network the smearing:
feed plaquettes built at the same cumulative APE levels the GEVP uses
($0,2,4,6$), stacked as input channels ($3$ spatial planes $\times$ 4 levels
$=12$ channels), with parallel transport still built from thin links.
Smearing is spatial-only, so the single-timeslice transfer-matrix condition
is untouched. The comparison also becomes maximally sharp: both methods now
see \emph{exactly the same} smearing levels --- the GEVP may take one global
linear combination of them, while GELT may mix them nonlinearly and
site-by-site. Any gap is attributable to the network.

The first attempt hit the scale-flat runaway of Section~\ref{sec:pitfalls}
(smooth coherent channels $\to$ constructive attention sums $\to$ squared
per-layer gain $\to$ $C(0)\sim10^{73}$ by epoch 2); with the
$(\log C(0))^2$ scale pin the warm-started retrain is stable throughout ---
$C(0)$ held at $\mathcal{O}(1)$, no overfit turn-up, validation flat on the
floor --- reaching best validation loss
\begin{equation}
  \mathcal{L}_{\rm val} = -0.6185
  \quad\Longleftrightarrow\quad
  m\,\at \approx 0.329 ,
\end{equation}
i.e.\ the Rayleigh quotient \emph{saturates the transfer-matrix bound at the
anchor mass}: the converged loss is the glueball mass, realized. On test
(Figure~\ref{fig:run5}):

\begin{table}[htbp]
  \centering
  \begin{tabular}{lcccc}
    \toprule
    operator & $\meff(1)$ & $\meff(2)$ & $\meff(3)$ & $\meff(4)$ \\
    \midrule
    GELT (smeared input)      & $\mathbf{0.370\pm0.016}$ & $\mathbf{0.333\pm0.020}$ & $0.345\pm0.029$ & --- \\
    classical GEVP            & $0.398\pm0.016$ & $0.357\pm0.025$ & $0.333\pm0.036$ & $0.284\pm0.050$ \\
    GEVP $+$ GELT (5 ops)     & $0.366\pm0.017$ & $0.333\pm0.020$ & $0.347\pm0.030$ & $0.319\pm0.048$ \\
    \bottomrule
  \end{tabular}
  \caption{Run 5 final results, blocked jackknife on the 400 test
  configurations. Units of $\at$. The learned single operator plateaus from
  $\Delta=1$; the enlarged five-operator GEVP collapses onto it.}
  \label{tab:run5}
\end{table}

Since GELT and the GEVP are measured on the same test configurations, their
error bars are correlated and comparing them independently would be
dishonest. The significance test is the blocked jackknife of the
\emph{difference} $\meff^{\rm GELT}-\meff^{\rm GEVP}$, in which the shared
fluctuations cancel:
\begin{equation}
  \Delta{=}1:\ -0.028\pm0.007\ (3.9\sigma), \qquad
  \Delta{=}2:\ -0.024\pm0.014\ (1.8\sigma), \qquad
  \Delta{=}3:\ +0.012\pm0.021 .
\end{equation}
The learned operator sits \emph{significantly below} the classical GEVP
exactly where excited-state contamination lives (small $\Delta$) and is
statistically identical on the plateau: a better operator measuring the same
physics --- which is precisely what the variational bound demands of a
legitimate win. Three corroborating observations:
\begin{itemize}
  \item \textbf{The win criterion is met.} GELT is flat at $\approx0.33$--$0.37$
  from $\Delta=1$ (in Figure~\ref{fig:run5} it hugs the anchor line through
  $\Delta\approx6$) while the GEVP is still descending at $\Delta=1$--$2$.
  Even GELT's most contaminated point, the $\Delta{=}0{\to}1$ Rayleigh mass
  $0.392$, sits below the GEVP's $\meff(1)$.
  \item \textbf{One learned operator subsumes the hand-built basis.} Offered
  the classical four-level ladder \emph{plus} GELT, the variational solver
  picks essentially pure GELT --- the enlarged GEVP equals GELT's own
  $\meff$ to within $\sim0.004$ at $\Delta=1$. Against the Run-4b null, this
  separation is cleanly attributable to the network's nonlinear, spatially
  resolved mixing of the same smearing levels.
  \item \textbf{The network adds content, not just smearing.} The ladder
  correlations of the trained operator are
  $|r| = 0.45, 0.76, 0.87, 0.91$ for levels $0,2,4,6$ --- it lives at the
  deep end of the ladder --- and training grew the ladder-\emph{independent}
  variance from $\sim8\%$ to $\sim18\%$ \emph{while the loss improved}: the
  network's own content, not residual noise, drives the win.
\end{itemize}

\begin{figure}[tbp]
  \centering
  \includegraphics[width=\textwidth]{glueball_gelt_GELT_beats_GEVP.png}
  \caption{Run 5 (final): GELT with multi-level smeared input channels
  ($[0,2,4,6]$). Left: with the scale pin the training is stable and the
  validation Rayleigh loss sits on the variational floor,
  $-0.6185 \leftrightarrow m\,\at\approx0.329$ (validation below training is
  an estimator artifact: the 6-configuration batch estimator is noise-biased
  toward zero, the 200-configuration validation estimate is cleaner). Right:
  test effective masses. The learned operator (green) is flat on the anchor
  (dashed) from $\Delta=1$, below the still-descending classical GEVP (red)
  at small $\Delta$; the enlarged GEVP $+$ GELT basis (purple) collapses onto
  pure GELT.}
  \label{fig:run5}
\end{figure}

\subsection{Quoting the result like a lattice paper: cosh fits and the
ground-state overlap $A_0$}
\label{sec:overlap}

Table~\ref{tab:run5} compares effective masses point by point. That is how
lattice results are \emph{plotted}, but not how they are \emph{quoted}: the
convention of the field is to quote the mass from a fit of $C(\Delta)$ to the
single-state form over a plateau window, drawn as a horizontal band through
the $\meff$ points, and to quote \emph{operator quality} --- which is what
the GELT-vs-GEVP comparison is actually about --- as the \emph{ground-state
overlap}, the number Morningstar and Peardon use to rank their smeared and
fuzzed operators \cite{MorningstarPeardon1999}. This subsection re-expresses
the Run-5 result in that standard form (Figure~\ref{fig:overlap}).

Each operator's connected correlator is fitted to the periodic
(finite-$N_t$) single-state form of Section~\ref{sec:correlators},
\begin{equation}
  C(\Delta) \;\approx\;
  A\,\big[\, e^{-m\,\at\,\Delta} + e^{-m\,\at\,(N_t-\Delta)} \,\big],
  \label{eq:coshfit}
\end{equation}
over one \emph{shared} window $\Delta\in[2,7]$ --- comparability beats
per-operator optimality, and starting at $\Delta=2$ keeps the GEVP's residual
$\Delta=1$ contamination out of \emph{its} mass fit: the contamination is
what the overlap should report, not what the fit should absorb. Extrapolating
the fitted ground-state piece back to $\Delta=0$ and comparing with the
measured $C(0)=\sum_{n>0}|\langle0|\hat{\bar O}|n\rangle|^2$ gives the
\textbf{ground-state overlap fraction}
\begin{equation}
  A_0 \;=\; \frac{A\,\big(1+e^{-m\,\at\,N_t}\big)}{C(0)} \;\in\; (0,1] ,
  \label{eq:overlap}
\end{equation}
the fraction of the operator's spectral weight carried by the glueball;
$A_0=1$ is a pure ground-state interpolator. This closes a conceptual loop:
the Rayleigh quotient \eqref{eq:rayleighquotient} is (up to the wrap term) a
weighted average $\sum_n w_n\,e^{-m_n\at}$ with weights summing to one and
$w_{\rm ground}=A_0$, so saturating the variational bound \emph{is} the
statement $A_0\to1$ --- the quantity quoted here is the quantity the
training maximizes. Objective and metric coincide.

The comparison also needs the right comparator. A single learned operator
should not be measured against a whole basis but against the best
\emph{single} operator the basis contains; the GEVP supplies exactly that
(Section~\ref{sec:gevp}): the ground-state eigenvector $v_0$ of
\eqref{eq:gevp}, solved once at reference times $(t_0,t_d)=(1,2)$, defines
the fixed-vector \emph{projected operator}
$\Obar_{\rm proj}(t)=\sum_i (v_0)_i\,\Obar_i(t)$ --- the variational optimum
in the span of the four smearing levels \cite{Blossier2009} --- whose scalar
correlator is fitted like any other. Errors follow the same discipline as
before: the $\sigma_\Delta$ entering the (diagonal) $\chi^2$ are fixed from
the full-sample blocked jackknife, the \emph{entire} fit --- including
$v_0$ --- is redone inside each of the 40 delete-block samples, and the
GELT$-$GEVP differences are jackknifed directly so the shared-ensemble
fluctuations cancel. On the 400 test configurations:

\begin{table}[htbp]
  \centering
  \begin{tabular}{lccc}
    \toprule
    operator & $m\,\at$ (fit) & $A_0$ & $\chi^2/\mathrm{dof}$ \\
    \midrule
    GELT (learned)    & $0.332\pm0.027$ & $\mathbf{0.903\pm0.047}$ & $0.05$ \\
    GEVP-projected    & $0.340\pm0.030$ & $0.837\pm0.056$          & $0.04$ \\
    APE$\times6$      & $0.340\pm0.032$ & $0.805\pm0.052$          & $0.11$ \\
    \bottomrule
  \end{tabular}
  \caption{Cosh fits \eqref{eq:coshfit} on the shared window
  $\Delta\in[2,7]$, blocked jackknife on the 400 test configurations
  ($\mathrm{dof}=4$). The $\chi^2/\mathrm{dof}$ values are artificially low
  because neighbouring $C(\Delta)$ points are strongly correlated while the
  $\chi^2$ is diagonal --- a fit-quality heuristic, not a goodness-of-fit
  test; the quoted errors come from the jackknife, which propagates the
  correlations correctly.}
  \label{tab:overlap}
\end{table}

\noindent with correlated differences
\begin{equation}
  \Delta(m\,\at) \;=\; -0.008\pm0.014\ \ (0.6\sigma), \qquad
  \Delta A_0 \;=\; +0.066\pm0.031\ \ (2.1\sigma).
  \label{eq:overlapdiff}
\end{equation}
The headline, in the field's own language: \emph{both operators measure the
same mass, but the learned operator carries $90\%$ of its spectral weight on
the ground state against $84\%$ for the optimal combination of the entire
classical basis.} Two readings of Table~\ref{tab:overlap} sharpen it:
\begin{itemize}
  \item \textbf{The GEVP barely improves on its best member; GELT does.} The
  optimal classical combination gains only $\sim0.03$ of overlap over
  APE$\times$6 alone ($0.837$ vs $0.805$), while GELT gains $\sim0.10$ over
  the full basis: one learned operator adds more ground-state overlap than
  the whole Morningstar--Peardon construction adds over its best single
  level --- the fit-language restatement of the enlarged-GEVP collapse of
  Section~\ref{sec:run5}.
  \item \textbf{$2.1\sigma$ and $3.9\sigma$ are complementary, not in
  tension.} The $3.9\sigma$ of Section~\ref{sec:run5} lives
  entirely at $\Delta=1$, where the contamination is largest; the shared fit
  window starts at $\Delta=2$, so $\Delta A_0$ never sees GELT's cleanest
  advantage and reconstructs it only by extrapolation. $A_0$ is the
  \emph{physical} operator-quality statement; the correlated $\meff$
  difference at $\Delta=1$ remains the sharpest single-point evidence of the
  same fact.
\end{itemize}

\begin{figure}[tbp]
  \centering
  \includegraphics[width=\textwidth]{glueball_overlap_run5.png}
  \caption{The Run-5 result in standard spectroscopic form. Left: test
  $\meff(\Delta)$ for GELT (green) and the projected GEVP operator (red)
  with their fitted-mass bands over the shared window $\Delta\in[2,7]$ ---
  the quoted mass is the band, not any single point; both bands sit on the
  classical anchor (dashed). Right: the ground-state fraction
  $\rho(\Delta)=[C(\Delta)/C(0)]\,/\,\cosh_{\mathrm{ref}}(\Delta)$, with
  $\cosh_{\mathrm{ref}}$ the fitted GELT ground state normalised to
  $\rho(0)=1$. A pure ground-state interpolator is flat at its $A_0$
  (dashed lines); excited-state contamination is the excess at small
  $\Delta$. Thin links collapse to $\rho\approx0.25$ immediately; APE$\times$6
  and the projected GEVP descend through $\Delta\approx3$ toward
  $A_0\approx0.80$--$0.84$; GELT is flat at $A_0\approx0.90$ from
  $\Delta=1$. (The classical curves' dip slightly below their $A_0$ lines
  near $\Delta\approx4$--$5$ is noise plus the cosh wrap, well within the
  error bars --- an exact spectral sum would approach $A_0$ from above.)}
  \label{fig:overlap}
\end{figure}

\subsection{Replication on an independent ensemble}
\label{sec:replication}

Everything so far is one ensemble. To close that caveat, the entire chain
--- heat-bath/overrelaxation sampling, from-scratch training, checkpoint
selection, test-split analysis --- was repeated with a different sampler
seed: a second, fully independent $12^3\times24$, $\beta=2.4$, $\xi=3$,
$N=2000$ ensemble, the same network and hyperparameters, no reuse of the
Run-5 weights. Nothing about the procedure was tuned on the new data.

The replication reproduces the operator-quality result --- and sharpens a
point about what does \emph{not} replicate. On the new ensemble every
operator, classical and learned alike, reads $\sim1\sigma$ heavier: the
classical GEVP sits flat at $\meff\approx0.38$--$0.40$ rather than
descending through $0.36$ toward $0.33$, and the fitted masses come out
$0.37$--$0.40$. Correspondingly, training converged to a best validation
Rayleigh loss of $-0.564$ rather than $-0.619$ --- which decodes to
$m\approx0.395$, i.e.\ the variational floor \emph{of this ensemble's own
correlator}, saturated to $\sim0.004$ exactly as before. The Rayleigh floor
is an empirical property of the ensemble, not a universal target: what the
loss saturates is the transfer-matrix bound as realized on the sampled
data, so the mass fluctuates between independent ensembles (here by
$0.04\pm0.04$, an ordinary $\sim1\sigma$ statistical difference on 2000
configurations --- the ``anchor'' carries its own ensemble error), while
the operator-quality statements are what should --- and do --- replicate.
On the new test split ($400$ configurations, same shared window
$\Delta\in[2,7]$):

\begin{table}[htbp]
  \centering
  \begin{tabular}{lccc}
    \toprule
    operator & $m\,\at$ (fit) & $A_0$ & $\chi^2/\mathrm{dof}$ \\
    \midrule
    GELT (learned)    & $0.374\pm0.031$ & $\mathbf{1.013\pm0.062}$ & $0.01$ \\
    GEVP-projected    & $0.400\pm0.038$ & $0.925\pm0.071$          & $0.08$ \\
    APE$\times6$      & $0.390\pm0.037$ & $0.901\pm0.068$          & $0.05$ \\
    \bottomrule
  \end{tabular}
  \caption{The replication ensemble's analogue of Table~\ref{tab:overlap}:
  cosh fits on the shared window $\Delta\in[2,7]$, blocked jackknife on the
  400 test configurations. $A_0$ exceeds $1$ only within its error ---
  consistent with a pure ground-state interpolator.}
  \label{tab:replication}
\end{table}

\noindent with correlated differences
\begin{equation}
  \Delta(m\,\at) \;=\; -0.025\pm0.013\ \ (1.9\sigma), \qquad
  \Delta A_0 \;=\; +0.089\pm0.030\ \ (2.9\sigma),
  \label{eq:replicationdiff}
\end{equation}
and a correlated $\meff$ difference of $-0.038\pm0.008$ ($4.5\sigma$) at
$\Delta=1$ and $-0.036\pm0.012$ ($2.9\sigma$) at $\Delta=2$ --- the Run-5
pattern, if anything stronger. The enlarged GEVP again collapses onto the
learned operator (the classical ladder adds $\sim0.005$ at $\Delta=1$,
within noise), and the smearing-ladder correlations are nearly identical
($r=0.36/0.75/0.89/0.93$ against APE$\times0/2/4/6$).

Since the two ensembles are statistically independent, the two overlap
differences combine: inverse-variance weighting of $+0.066\pm0.031$ (Run 5)
and $+0.089\pm0.030$ (replication) gives
\begin{equation}
  \Delta A_0^{\rm combined} \;=\; +0.078\pm0.022 \qquad (3.6\sigma).
  \label{eq:combined}
\end{equation}
The learned operator's excess ground-state overlap over the optimal
classical combination is not a feature of one ensemble.

\begin{figure}[tbp]
  \centering
  \includegraphics[width=\textwidth]{glueball_overlap_ens1.png}
  \caption{Figure~\ref{fig:overlap} repeated on the independent replication
  ensemble: $\meff(\Delta)$ with fitted-mass bands (left) and the
  ground-state fraction $\rho(\Delta)$ (right). The overall mass scale sits
  $\sim1\sigma$ higher than on the original ensemble --- both operators
  agree on that shift, as they must --- while the operator-quality ordering
  (GELT flat at its $A_0$ from small $\Delta$, the classical operators
  descending into it) is unchanged.}
  \label{fig:replication}
\end{figure}

% ==============================================================================
\section{What the learned operator attends to}
\label{sec:attention}
% ==============================================================================

The validation above establishes \emph{that} the learned operator is better;
it says nothing about \emph{how}. Because every GELT attention score is a
gauge invariant --- $\mathrm{Re}\,\Tr[Q^\dagger\tilde K]$ between parallel
transported adjoint fields --- the attention weights $\alpha_n(x)$ are not an
implementation detail but a measurable field: a per-site probability
distribution over the $24$ neighbours in the $L_1$ ball of radius $R=2$, plus
the site itself. This section reads that distribution out of the Run-5
checkpoint and asks what structure the optimal operator has learned. It is a
first look, on deliberately modest statistics, not a finished study.

\paragraph{Setup.} A fresh $12^3\times24$, $\beta=2.4$, $\xi=3$ ensemble of
$N=32$ configurations was sampled with a new seed, so it is held out by
construction; the Run-5 checkpoint was loaded unchanged and run in inference
over all $32\times24=768$ timeslices. Three statistics are accumulated per
(layer, head): the \emph{attention range}
$\ell_{\mathrm{att}}(x)=\sum_n\alpha_n(x)\,|\Delta x_n|_1$, the offset entropy
$-\sum_n\alpha_n\log\alpha_n$ (maximal value $\log 25=3.22$), and a per-site
axial anisotropy --- the relative spread, across the three spatial axes, of the
attention mass at $|\Delta x|_1=1$ (maximal value $\sqrt2=1.41$, attained when
all of it sits on one axis). Two further arms follow the rule that a
correlation is not an explanation: the per-timeslice Spearman correlation of
$\ell_{\mathrm{att}}(x)$ against the smeared spatial-plaquette density $o(x)$
(the physical field this operator measures --- large $o$ means \emph{small}
Wilson action density, since $S\propto\sum_x(1-o(x)/3)$), and a head-ablation
intervention in which head $h$'s column of layer $l$'s output mixing matrix is
zeroed and the degradation in the Rayleigh loss measured. Errors on both
validation arms respect the correct independent unit: the ablation
$\Delta$loss is a delete-one jackknife over configurations of the
\emph{correlated} difference (ablated and intact operators evaluated on the
same configurations, so the shared ensemble noise cancels), and the Spearman
error is the spread of per-configuration means, not of per-slice values ---
slices within a configuration are correlated over $\sim1/m$, which is the very
quantity this report measures. The whole study runs in about an hour on a
laptop CPU.

\paragraph{The checkpoint transfers.} On this previously unseen chain the
operator returns $C(1)/C(0)=0.648$, i.e.\ a Rayleigh loss of $-0.538$ and
$m\,\at\approx0.43$. This is a third independent realization of the
ensemble-specific variational floor discussed in
Section~\ref{sec:replication} ($0.33$ on Run 5, $0.395$ on the replication
ensemble), and it confirms that the trained weights are not tuned to their
training chain. The number should be read with its statistical width in mind:
the same quantity on the first $8$ configurations of this same chain reads
$m\,\at\approx0.375$, so a Rayleigh ratio on tens of configurations carries an
uncertainty of order $0.05$ --- which is exactly why the intervention arm
below needed error bars before it could be interpreted at all.

\begin{table}[htbp]
  \centering
  \begin{tabular}{ccccccc}
    \toprule
    layer & head & $\ell_{\mathrm{att}}$ & entropy & anisotropy
          & axis split & Spearman$(\ell,o)$ \\
    \midrule
    1 & 0 & $1.69\pm0.15$ & $2.24$ & $0.56$ & $0.40/0.31/0.29$ & $-0.139(1)$ \\
    1 & 1 & $1.87\pm0.06$ & $1.04$ & $0.99$ & $0.00/1.00/0.00$ & $-0.009(1)$ \\
    2 & 0 & $1.93\pm0.05$ & $0.35$ & $1.37$ & $0.00/1.00/0.00$ & $-0.086(1)$ \\
    2 & 1 & $2.00\pm0.00$ & $0.004$ & $1.41$ & $1.00/0.00/0.00$ & $+0.397(2)$ \\
    3 & 0 & $2.00\pm0.01$ & $0.07$ & $1.41$ & $0.00/1.00/0.00$ & $-0.088(1)$ \\
    3 & 1 & $\mathbf{0.30\pm0.18}$ & $0.89$ & $0.79$ & $0.81/0.00/0.19$
      & $-0.328(2)$ \\
    4 & 0 & $1.59\pm0.21$ & $1.25$ & $1.36$ & $0.99/0.00/0.00$ & $-0.230(2)$ \\
    4 & 1 & $1.44\pm0.12$ & $2.22$ & $0.75$ & $0.00/0.64/0.36$
      & $\mathbf{-0.564(2)}$ \\
    \bottomrule
  \end{tabular}

  \vspace{1.2em}

  \begin{tabular}{cccc}
    \toprule
    layer & head & $\Delta$loss (ablate) & significance \\
    \midrule
    1 & 0 & $-0.0023\pm0.0010$ & $-2.3\sigma$ \\
    1 & 1 & $-0.0002\pm0.0005$ & $-0.4\sigma$ \\
    2 & 0 & $-0.0025\pm0.0011$ & $-2.2\sigma$ \\
    2 & 1 & $-0.0003\pm0.0017$ & $-0.2\sigma$ \\
    3 & 0 & $-0.0010\pm0.0014$ & $-0.7\sigma$ \\
    3 & 1 & $\mathbf{+0.0787\pm0.0145}$ & $\mathbf{+5.4\sigma}$ \\
    4 & 0 & $+0.0005\pm0.0055$ & $+0.1\sigma$ \\
    4 & 1 & $+0.0058\pm0.0036$ & $+1.6\sigma$ \\
    \bottomrule
  \end{tabular}
  \caption{Attention statistics (top) and head-ablation intervention (bottom)
  of the Run-5 operator on the held-out seed-7 ensemble ($32$ configurations,
  $768$ timeslices, $1728$ sites each). The $\ell_{\mathrm{att}}$ error is the
  \emph{site-to-site} spread, not a mean error: it is the
  configuration-dependent part of the attention, the part a fixed convolution
  kernel could not supply. ``Axis split'' is the fraction of sites at which
  the head selects each spatial axis for its $|\Delta x|_1=1$ attention mass.
  Spearman errors are configuration-level. Baseline Rayleigh loss $-0.538$;
  a positive $\Delta$loss means removing the head makes the operator worse,
  and the error is a delete-one jackknife over configurations of the
  correlated difference.}
  \label{tab:attention}
\end{table}

\paragraph{One head does the work.} The ablation is strikingly concentrated:
head 1 of layer 3 costs $+0.0787\pm0.0145$ in Rayleigh loss when removed ---
$5.4\sigma$, an order of magnitude more than any other head, and roughly a
seventh of the entire signal --- while no other head reaches $2\sigma$ in the
positive direction. That same head is the only \emph{local} one in the
network: $\ell_{\mathrm{att}}=0.30$ with $71\%$ of its attention mass on the
site itself, against six heads that sit at or near the maximal range
$\ell_{\mathrm{att}}=2$. The optimal operator is therefore not a democratic
ensemble of scales but one content-dependent local head embedded in a stack of
long-range ones.

Two heads carry a mildly \emph{negative} $\Delta$loss (layer 1 head 0 and
layer 2 head 0, both $\approx-0.0024$ at $2.2$--$2.3\sigma$): removing them
marginally improves the Rayleigh ratio on this ensemble. Eight cells were
tested, so $\sim0.2$ excursions beyond $2.2\sigma$ are expected by chance and
two is only a mild excess; but the direction is also what mild overfitting
would produce, since these weights were optimized on the Rayleigh ratio of a
\emph{different} ensemble and are being scored here on a fresh one. Either way
the effect is $30\times$ smaller than the layer-3 head and does not bear on
the operator's measured quality.

\paragraph{No emergent coarsening.} The interpretability program's
Study-1 hypothesis --- that a learned operator should widen its receptive
field with depth, an RG-like coarsening no fixed architecture imposes --- is
\emph{not} supported here. $\ell_{\mathrm{att}}$ is essentially flat across
layers at the ceiling $\ell=2$, and the one departure (layer 3, head 1) goes
the wrong way. Worse, two heads are effectively degenerate: layer 2 head 1 has
entropy $0.004$ and $\ell_{\mathrm{att}}=2.000\pm0.000$, and layer 3 head 0 is
close behind, meaning they place all of their mass deterministically on the
outer $|\Delta x|_1=2$ shell regardless of the configuration. Those heads are
fixed kernels, not measurements: whatever they contribute, a convolution could
have supplied it. This is a negative result about the current
architecture --- with $R=2$ the ball is too small for a range hierarchy to have
anywhere to go --- and it is worth reporting as such, since it identifies $R$
and the number of layers, not the attention mechanism, as what to vary next.

\paragraph{Attention reaches further where the action is dense.} The
ground-truth arm is the more encouraging one. Layer 4 head 1 --- the highest
entropy head in the network, i.e.\ the most genuinely distributed --- has
$\ell_{\mathrm{att}}(x)$ anticorrelating with the smeared plaquette density at
Spearman $-0.56$. Since large $o(x)$ is \emph{small} action density, the sign
says the head extends its reach precisely on the disordered, high-action
regions and contracts on smooth vacuum. That is the behaviour one would want
of a glueball operator, and it is measured against a physical field rather
than asserted from a picture. Layer 3 head 1, the load-bearing local head,
shows the same sign more weakly ($-0.33$).

\paragraph{The axial anisotropy is a globally fixed axis.} Four heads sit at
anisotropy $\gtrsim1.36$ out of a maximum of $1.41$: at each site, their
distance-one attention is concentrated on a single spatial axis. The $0^{++}$
channel is a cubic-group scalar, and nothing in GELT enforces rotational
symmetry (only gauge equivariance; the shortest-path transport average is
rotation-symmetric, the learned weights are not). The question is whether the
selected axis varies from site to site, in which case isotropy is restored in
the zero-momentum sum, or is globally fixed, in which case the operator is not
a cubic-group scalar. The axis-split column settles it, and the answer is the
unfavourable one: layer 1 head 1, layer 2 head 0 and layer 3 head 0 select the
same axis at $100.0\%$ of all $1.3$ million site--slice samples, layer 2
head 1 selects a different single axis at $100.0\%$, and layer 4 head 0 at
$99\%$. These are not content-driven choices; they are frozen directions.

Two observations keep this from being fatal. First, the axis-locked heads are
exactly the low-entropy, maximal-range, ablation-irrelevant ones --- the
degenerate kernels of the previous paragraph --- while the two heads that
matter physically are the two most isotropic ones: the load-bearing layer-3
head splits $0.81/0.00/0.19$ and the best-correlated layer-4 head splits
$0.00/0.64/0.36$. Second, a spin-2 admixture in the interpolating operator
contaminates the correlator with \emph{heavier} states, biasing an extracted
mass upward, not downward; it cannot manufacture the ground-state overlap
excess of Section~\ref{sec:overlap}. What it does mean is that the learned
operator is a $0^{++}$ interpolator only approximately, and that explicitly
projecting onto the cubic $A_1^{++}$ representation --- averaging the operator
over the $24$ lattice rotations, or symmetrizing the architecture --- is a
concrete and likely profitable next step.

\begin{figure}[tbp]
  \centering
  \includegraphics[width=\textwidth]{glueball_attention.png}
  \caption{Attention of the Run-5 operator as a measurement. Top:
  $\ell_{\mathrm{att}}$ versus depth with site-to-site spread as error bars
  (left), the config-averaged radial profile (centre), and the head-ablation
  matrix (right). Bottom: the distribution of the per-site range, whose width
  is the configuration dependence (left), and one $z=0$ slice of
  $\delta\ell_{\mathrm{att}}(x)$ next to the smeared density $\delta o(x)$ for
  the most strongly correlated head (centre, right).}
  \label{fig:attention}
\end{figure}

\begin{figure}[tbp]
  \centering
  \includegraphics[width=\textwidth]{glueball_attention_kernels.png}
  \caption{Config-averaged attention kernels $\bar\alpha(\Delta x)$ on the
  $\Delta x_3=0$ plane, per (layer, head), on a logarithmic scale. This is
  the \emph{static} part of the attention only --- the part expressible as a
  fixed convolution; the claim of Table~\ref{tab:attention} lives in the
  site-to-site residual around these means.}
  \label{fig:attentionkernels}
\end{figure}

\paragraph{Attention range versus correlation length: a measurement this
lattice does not permit.} The natural companion to the localization question
is whether the learned attention range tracks a physical length --- whether
$\ell_{\mathrm{att}}$ grows with the correlation length $\xi$ as the coupling
is varied. The relevant length is the \emph{spatial} one, in the units the
attention offsets are counted in:
\begin{equation}
  \xi_s \;=\; \frac{1}{m\,\as} \;=\; \frac{1}{\xi\, m\,\at},
  \label{eq:xis}
\end{equation}
since $\ell_{\mathrm{att}}$ is built from spatial displacements $|\Delta x|_1$
measured in $\as$. Before training one operator per coupling, it is worth
asking whether $\xi_s$ moves at all. It is a purely classical question --- one
ensemble and one GEVP per $\beta$, no network --- and the answer decides
whether the study is possible:

\begin{table}[htbp]
  \centering
  \begin{tabular}{ccccc}
    \toprule
    $\beta$ & $\beta_s=\beta/\xi$ & $m\,\at$ (GEVP, $\Delta=2$) & $m\,\as$
            & $\xi_s$ \\
    \midrule
    2.1 & 0.70 & $0.5812\pm0.0246$ & 1.744 & $0.57\pm0.02$ \\
    2.3 & 0.77 & $0.4009\pm0.0128$ & 1.203 & $0.83\pm0.03$ \\
    2.4 & 0.80 & $0.3330\pm0.0106$ & 0.999 & $1.00\pm0.03$ \\
    2.5 & 0.83 & $0.3026\pm0.0105$ & 0.908 & $1.10\pm0.04$ \\
    2.7 & 0.90 & $0.3325\pm0.0113$ & 0.998 & $1.00\pm0.03$ \\
    \bottomrule
  \end{tabular}
  \caption{Classical $\beta$-scan on the $12^3\times24$, $\xi=3$ lattice
  family ($N=2000$ per coupling, multi-level GEVP, blocked jackknife). The
  $\beta=2.4$ row is the anchor ensemble of Section~\ref{sec:runs} and
  reproduces its established value $0.333\pm0.011$, which validates the scan.
  $\xi_s$ is Eq.~\eqref{eq:xis}, in units of the spatial lattice spacing.}
  \label{tab:betascan}
\end{table}

\noindent The decisive number is not the span but the magnitude: $\xi_s$ never
exceeds $1.1$ lattice spacings anywhere in the window. The attention offsets
are \emph{integers}, $|\Delta x|_1\in\{0,1,2,\dots\}$, so the entire variation
of the quantity $\ell_{\mathrm{att}}$ would have to track lives inside the
first offset shell. A factor of two in a sub-lattice length is not a factor of
two in anything the network can represent, and training one operator per
$\beta$ would return a flat scatter whose flatness measured the lattice
spacing rather than the attention.

It is worth being precise about what is and is not small here, because the two
are easily confused. The \emph{box} is large: $L/\xi_s \approx 12$ correlation
lengths across, which is the regime in which finite-volume effects are
absent, not present. What is coarse is the \emph{spacing}. At $\xi=3$ the
spatial coupling is $\beta_s=\beta/\xi\approx0.8$ --- strong coupling --- so
$\as$ is large and the entire physical extent of the glueball falls inside a
single grid cell. The lattice has plenty of cells; the object occupies about
one of them. Refining $\as$, not enlarging $L$, is what the measurement would
need.

One caveat on the high-$\beta$ end. $m\,\at$ falls monotonically from
$\beta=2.1$ to $\beta=2.5$ and then \emph{rises} at $\beta=2.7$, where
continuum scaling demands it keep falling. With $L/\xi_s\approx12$ this cannot
be a finite-volume squeeze. The natural reading is excited-state
contamination at fixed $\Delta$: the mass is quoted at $\Delta=2$, and since
$\at$ shrinks as $\beta$ grows, that same $\Delta$ is a shorter \emph{physical}
Euclidean time in which to filter excited states, biasing $\meff$ upward. A
second systematic points the same way --- the conversion $m\,\as=\xi\,m\,\at$
uses the \emph{bare} anisotropy, while the renormalized one differs and drifts
with $\beta$. Both mean the $\beta=2.7$ entry should be read as an upper bound
on the mass, hence a \emph{lower} bound on $\xi_s$: the true window may be
somewhat wider than the tabulated factor $1.9$. It would have to be very much
wider --- $\xi_s$ clearing several spacings --- to change the conclusion, but
the honest statement is that the high-$\beta$ end is bounded rather than
measured.

The honest conclusion is that this is a statement about the lattice, not about
the architecture. Reaching $\xi_s\gg1$ requires either a larger volume at
larger $\beta$ (cost $\propto L^3$) or a milder anisotropy: since
$\beta_s=\beta/\xi$ sets $\as$, lowering $\xi$ from $3$ towards $1.5$ at fixed
$\beta$ would double $\beta_s$ and should push $\xi_s$ comfortably past one
spacing while keeping $m\,\at$ resolvable in time. Both are ensemble campaigns
rather than analyses, and neither is attempted here. What the scan does
deliver, at the cost of one unattended run rather than four trainings, is a
quantitative reason for the omission --- and, incidentally, an independent
reproduction of the anchor mass.

\paragraph{What this section is not.} The statistics above are honest but
small, and it is worth recording what taking the independent unit seriously
cost. The ablation was first run on $8$ configurations with no error bars,
where the layer-3 head read $+0.127$ and four of the eight cells came out
negative. At $32$ configurations with a correlated jackknife the same head
reads $+0.0787\pm0.0145$ and the negative cells collapse to the
$\lesssim0.003$ level: the point estimate moved by more than $3\sigma$ of its
own final error. A Rayleigh ratio on a handful of configurations is simply not
a measurement --- the same caution applies to the $m\,\at\approx0.43$ quoted
above, whose $8$-configuration value on the same chain was $0.375$. The
Spearman correction was milder than the ablation's: configuration-level errors
come out only $1.5\times$ the naive slice-level ones, so the strong
correlations ($-0.56$, $-0.33$) were never in doubt, but the quoted precision
now means what it says.

What remains open is the receptive-field question, and no amount of reanalysis
can settle it: deciding whether the absence of depth-wise coarsening is a
property of the attention mechanism or merely of $R=2$ requires
\emph{retraining} at larger $R$ on the full $2000$-configuration ensemble,
which is GPU work. Until then the negative result stands only for the network
as trained. Note that the $\beta$-scan above makes the same point from the
physics side: with $\xi_s\lesssim1$, a larger $R$ would give the attention
more room without giving it more to see.

% ==============================================================================
\section{Conclusions, caveats, outlook}
% ==============================================================================

The classical chain --- Wilson action, heat-bath/overrelaxation ensemble,
smeared operators, multi-level GEVP, anisotropic lattice, jackknife ---
resolves the $SU(2)$ $0^{++}$ ground state at $m\,\at\approx0.33$ on a
$12^3\times24$, $\beta=2.4$, $\xi=3$ ensemble, after the isotropic lattice
demonstrably could not. On top of it, a gauge-equivariant network trained on
the Rayleigh-quotient loss, constrained to a single timeslice so the
transfer-matrix bound applies, and fed the same smearing levels as the
classical basis, saturates the variational bound at the anchor mass and beats
the classical GEVP where it can legitimately be beaten: $3.9\sigma$ less
excited-state contamination at $\Delta=1$, identical mass on the plateau, and
a single learned operator that variationally subsumes the four-operator
hand-built basis. In the field's standard quoting (Section~\ref{sec:overlap}):
fitted masses $0.332(27)$ vs $0.340(30)$ --- the same physics --- with
ground-state overlap $A_0=0.90(5)$ against $0.84(6)$ for the optimal
classical combination, a correlated difference of $+0.066\pm0.031$; the
projected GEVP gains only $\sim0.03$ of overlap over its best single smearing
level, the learned operator $\sim0.10$ over the whole basis. A from-scratch
replication on an independently sampled ensemble
(Section~\ref{sec:replication}) reproduces all of the operator-quality
statements --- $\Delta A_0=+0.089\pm0.030$, a $4.5\sigma$ correlated $\meff$
difference at $\Delta=1$ --- and the two ensembles combine to
$\Delta A_0=+0.078\pm0.022$ ($3.6\sigma$).

Honest boundaries of the claim. (i) Two independent ensembles now, but still
one $(\beta,\xi)$ and coarse $\as$ --- no continuum statement is made, for
the learned operator exactly as for the classical one
(Section~\ref{sec:caveat}); and the replication shows the mass itself
fluctuates by $\sim1\sigma$ between ensembles, so the ``anchor'' carries its
own ensemble error. (ii) The $3.9\sigma$ and the overlap
differences are statements about \emph{operator quality}, not about the mass
value, on which the two methods agree --- as the variational principle says
they must.
(iii) Still open: a matched-parameter L-CNN baseline on the identical
per-timeslice task.
(iv) The attention readout of Section~\ref{sec:attention} rests on $32$
configurations, enough to resolve its one large intervention effect at
$5.4\sigma$ but not the sub-$0.01$ structure around it; and it reveals that
the learned operator is only approximately a cubic-group scalar.
(v) Whether the attention range tracks the correlation length is
\emph{not} measurable on this lattice family: the classical $\beta$-scan of
Table~\ref{tab:betascan} puts $\xi_s$ below $1.1$ spatial lattice spacings
everywhere it can be measured, which is smaller than the grid the attention
lives on.

Beyond validation, the trained operator opens a door the classical one
cannot: its gauge-invariant attention maps are a direct image of the spatial
loop structure the \emph{optimal} glueball operator attends to --- what the
learned $18\%$ beyond the smearing ladder actually computes is now an
interpretability question with a physical answer. A first readout
(Section~\ref{sec:attention}) already returns four concrete findings: the
operator's advantage is carried almost entirely by a single \emph{local}
attention head ($5.4\sigma$ under ablation, against nothing else above
$2\sigma$); the hypothesized RG-like widening of the receptive field with
depth does \emph{not} occur at $R=2$, several heads having collapsed to fixed
maximal-range kernels; the most distributed head measurably extends its reach
on high-action regions and contracts on smooth vacuum; and those frozen heads
are locked to single spatial axes, so the learned operator is a cubic-group
scalar only approximately. The last of these is the most actionable: an
explicit $A_1^{++}$ projection --- averaging over the $24$ lattice rotations
--- is cheap, physically motivated, and untested.

% ==============================================================================
\begin{thebibliography}{99}

\bibitem{Wilson1974}
K.~G. Wilson,
\emph{Confinement of quarks},
Phys.\ Rev.\ D \textbf{10} (1974) 2445.

\bibitem{Berg1980}
B.~Berg,
\emph{Plaquette-plaquette correlations in the SU(2) lattice gauge theory},
Phys.\ Lett.\ B \textbf{97} (1980) 401.

\bibitem{OsterwalderSeiler1978}
K.~Osterwalder and E.~Seiler,
\emph{Gauge field theories on a lattice},
Ann.\ Phys.\ \textbf{110} (1978) 440.

\bibitem{Creutz1980}
M.~Creutz,
\emph{Monte Carlo study of quantized SU(2) gauge theory},
Phys.\ Rev.\ D \textbf{21} (1980) 2308.

\bibitem{KennedyPendleton1985}
A.~D. Kennedy and B.~J. Pendleton,
\emph{Improved heatbath method for Monte Carlo calculations in lattice gauge
theories},
Phys.\ Lett.\ B \textbf{156} (1985) 393.

\bibitem{Creutz1987}
M.~Creutz,
\emph{Overrelaxation and Monte Carlo simulation},
Phys.\ Rev.\ D \textbf{36} (1987) 515.

\bibitem{APE1987}
M.~Albanese et al.\ [APE Collaboration],
\emph{Glueball masses and string tension in lattice QCD},
Phys.\ Lett.\ B \textbf{192} (1987) 163.

\bibitem{Lepage1989}
G.~P. Lepage,
\emph{The analysis of algorithms for lattice field theory},
in \emph{From Actions to Answers} (TASI 1989), World Scientific (1990).

\bibitem{MichaelTeasdale1983}
C.~Michael and I.~Teasdale,
\emph{Extracting glueball masses from lattice QCD},
Nucl.\ Phys.\ B \textbf{215} (1983) 433.

\bibitem{LuscherWolff1990}
M.~L\"uscher and U.~Wolff,
\emph{How to calculate the elastic scattering matrix in two-dimensional
quantum field theories by numerical simulation},
Nucl.\ Phys.\ B \textbf{339} (1990) 222.

\bibitem{Blossier2009}
B.~Blossier, M.~Della Morte, G.~von Hippel, T.~Mendes and R.~Sommer,
\emph{On the generalized eigenvalue method for energies and matrix elements
in lattice field theory},
JHEP \textbf{04} (2009) 094, arXiv:0902.1265.

\bibitem{MorningstarPeardon1999}
C.~J. Morningstar and M.~Peardon,
\emph{The glueball spectrum from an anisotropic lattice study},
Phys.\ Rev.\ D \textbf{60} (1999) 034509, arXiv:hep-lat/9901004.

\bibitem{MadrasSokal1988}
N.~Madras and A.~D. Sokal,
\emph{The pivot algorithm: A highly efficient Monte Carlo method for the
self-avoiding walk},
J.\ Stat.\ Phys.\ \textbf{50} (1988) 109.

\bibitem{Favoni2020}
M.~Favoni, A.~Ipp, D.~I. M\"uller and D.~Schuh,
\emph{Lattice gauge equivariant convolutional neural networks},
Phys.\ Rev.\ Lett.\ \textbf{128} (2022) 032003, arXiv:2012.12901.

\end{thebibliography}

\end{document}
